\section{Experiments}

\subsection{Dataset}
The experiments use simulated 3D diffraction data and corresponding real-space
complex volumes. The training split contains \todo{number} samples and the
validation/test split contains \todo{number} samples. Each sample has spatial
size $64^3$.

\subsection{Training Setup}
Unless otherwise specified, models are trained from scratch with Adam and an
initial learning rate of \todo{learning rate}. The default loss is far-field L1.
For small-sample debugging, we also use an overfitting protocol in which a fixed
subset of samples is used for both training and validation.

The small-sample protocol is used to answer a narrower question than general
validation accuracy: can the model fit a known fixed subset and produce
physically plausible object-domain amplitude and phase for those same samples?
This prevents generalization error from being confused with an optimization or
architecture-prior problem.

\subsection{Experiment Design}
The paper focuses on a small number of experiments that directly support the
main claim, rather than reporting every debugging run. The core comparison is:
\begin{itemize}
  \item AutoPhaseNN reproduction from scratch with the paper-style far-field
  objective.
  \item UMamba backbone replacement with the same data interface,
  post-processing, objective, and evaluation protocol.
  \item One representative diagnostic ablation, selected only if it helps
  explain why Fourier-domain fitting and real-space correctness diverge.
\end{itemize}
Learning-rate sweeps, checkpoint debugging, gradient probes, and implementation
sanity checks are treated as internal controls and are not reported as main
experiments unless they change the scientific conclusion.

\subsection{Evaluation Metrics}
We report fixed metric groups:
\begin{itemize}
  \item FT metrics: L1, MSE, RMSE, and relative L1 of diffraction modulus.
  \item Amplitude metrics: full-volume L1, MSE, RMSE, and relative L1.
  \item Phase metrics: wrapped phase error on the true support.
  \item Support metrics: L1, IoU, Dice score, and predicted support fraction.
\end{itemize}

\subsection{Implementation Details}
The AutoPhaseNN and UMamba pipelines share the same dataset interface,
post-processing logic, checkpoint format, and evaluation scripts wherever
possible. This controls for non-model factors in the comparison.
